\documentclass[12pt]{article}

% \usepackage{sbc-template}
\usepackage{graphicx,url}
\usepackage[utf8]{inputenc}
\usepackage[brazil]{babel}
\usepackage{amsmath} % Adicionado para melhor formatação matemática
\usepackage{algorithm}
\usepackage{algpseudocode}

% Tradução dos termos do algoritmo para Português
\renewcommand{\algorithmicrequire}{\textbf{Entrada:}}
\renewcommand{\algorithmicensure}{\textbf{Saída:}}
\algrenewcommand\algorithmicfor{\textbf{para}}
\algrenewcommand\algorithmicdo{\textbf{faça}}
\algrenewcommand\algorithmicif{\textbf{se}}
\algrenewcommand\algorithmicthen{\textbf{então}}
\algrenewcommand\algorithmicend{\textbf{fim}}
\algrenewcommand\algorithmicreturn{\textbf{retorna}}


\sloppy

\title{Avaliação Assintótica de Métodos de Estimação de Dimensão Fractal: Box Counting vs. Gliding Box}

\author{Seu Nome\inst{1}, Nome do Coautor\inst{1}}

\address{Departamento de Computação -- Instituição (Sigla)\\
  Cidade -- Estado -- Brasil
  \email{\{seuemail, email\}@instituicao.br}
}

\begin{document} 

\maketitle

\begin{abstract}
  This document presents a theoretical foundation for the asymptotic evaluation of fractal dimension estimation methods, specifically comparing the Box Counting and Gliding Box algorithms. Based on recent literature, it discusses the computational complexity (time and space) and statistical reliability of both approaches when applied to large datasets and complex spatial structures.
\end{abstract}
     
\begin{resumo} 
  Este documento apresenta a fundamentação teórica para a avaliação assintótica de métodos de estimação de dimensão fractal, comparando especificamente os algoritmos de Box Counting e Gliding Box. Com base na literatura recente, discute-se a complexidade computacional (tempo e espaço) e a confiabilidade estatística de ambas as abordagens quando aplicadas a grandes conjuntos de dados e estruturas espaciais complexas.
\end{resumo}

\section{Introdução}

A geometria e os conceitos fractais têm sido crescentemente aplicados em diversos campos da ciência para descrever a complexidade e a autossimilaridade na natureza \cite{saa2007}. A caracterização de estruturas irregulares frequentemente exige a estimação da sua dimensão fractal matemática, o que desencadeou o desenvolvimento de múltiplos algoritmos computacionais \cite{gonzato1998}. 

Entre as técnicas mais populares encontra-se o método de \textit{Box Counting} (Contagem de Caixas), que estima a dimensão baseando-se na cobertura do objeto por uma grade de caixas mutuamente exclusivas \cite{saa2007}. Apesar de sua simplicidade e ausência de restrições quanto à conectividade do objeto medido \cite{gonzato1998}, o método clássico esbarra em limitações estatísticas e problemas de resolução quando aplicado a grades finitas \cite{saa2007}. 

Para contornar esses problemas, o método \textit{Gliding Box} (Caixa Deslizante) foi proposto. Em vez de utilizar partições estáticas, este método desliza uma caixa de tamanho pré-definido por todo o plano de suporte, gerando amostras sobrepostas \cite{saa2007}. Embora isso melhore substancialmente a robustez estatística dos resultados \cite{saa2007}, a sobreposição introduz novos desafios computacionais.

O objetivo deste trabalho é apresentar a fundamentação teórica e a avaliação assintótica (complexidade de tempo e espaço) desses dois algoritmos, evidenciando as trocas (*trade-offs*) entre precisão estatística e viabilidade computacional no processamento de grandes volumes de dados.

\section{Fundamentação Teórica}

A análise fractal computacional exige a tradução de limites teóricos contínuos para algoritmos discretos. As subseções a seguir detalham os mecanismos dos métodos \textit{Box Counting} e \textit{Gliding Box}, bem como suas respectivas complexidades assintóticas.

\subsection{O Método Box Counting (BC)}

A dimensão de capacidade (frequentemente chamada de dimensão de *box-counting*) é formalmente definida pelo limite de $d_B = \lim_{\epsilon \to 0} \frac{\log N_B(\epsilon)}{\log(1/\epsilon)}$ \cite{liebovitch1989}, onde $N_B(\epsilon)$ é o número mínimo de caixas de tamanho $\epsilon$ necessárias para cobrir o conjunto. Na prática, este cálculo é realizado avaliando a relação linear em um gráfico log-log entre o número de caixas "cheias" e o tamanho da caixa \cite{gonzato1998}.

Uma implementação ingênua deste método pode ser ineficiente. Contudo, Liebovitch e Toth (1989) propuseram um algoritmo eficiente que codifica as coordenadas dos pontos em caixas usando operações binárias e os ordena \cite{liebovitch1989}. Essa implementação otimizada requer um tempo computacional da ordem de $O(N \log N)$ (limitado principalmente pelo tempo de ordenação) e um espaço de memória proporcional a $N d_e$, onde $N$ é o número de pontos de dados e $d_e$ é a dimensão de imersão (*embedding dimension*) \cite{liebovitch1989}. Outras abordagens focam na manipulação de telas virtuais (*virtual screens*) de baixo consumo de memória para estender a resolução sem sobrecarga no processamento \cite{gonzato1998}.

A principal restrição estatística do método BC está nos tamanhos maiores de caixas: quanto maior a caixa, menor o número de amostras $n(\epsilon)$, o que prejudica a representação estatística da distribuição da medida sobre o plano \cite{saa2007}.

\subsection{O Método Gliding Box (GB)}

Originalmente introduzido para a análise de lacunaridade, o método GB constrói as amostras deslizando uma caixa de área $\epsilon$ sobre o mapa de grade em todas as direções possíveis \cite{saa2007}. Como essa partição cria sobreposições, a medida definida nessas caixas não é estatisticamente independente \cite{saa2007}. 

No entanto, a grande vantagem do método GB reside na obtenção de um tamanho de amostra consideravelmente maior para qualquer escala, o que geralmente leva a estimativas estatísticas mais robustas, com menores erros padrão e redução do viés associado ao formato geral dos dados \cite{saa2007}.

\subsection{Avaliação Assintótica e Complexidade Computacional}

O incremento na confiabilidade estatística do GB vem com um custo computacional substancial. O algoritmo original do Gliding Box opera com uma complexidade de tempo de $O(N^D)$ para contar massas em todas as possíveis posições, e sua complexidade de espaço também figura como $O(N^D)$, onde $N$ é o tamanho da imagem em cada dimensão e $D$ é o número de dimensões \cite{barna2024}. Isso o torna proibitivo para conjuntos de dados massivos.

Recentemente, métodos otimizados para operações equivalentes ao Gliding Box foram sistematizados para mitigar essa complexidade \cite{barna2024}:
\begin{itemize}
    \item \textbf{Método \textit{Sides}:} Utiliza a sobreposição consecutiva para calcular a massa avaliando apenas as bordas que não se sobrepõem \cite{barna2024}. Mantém a complexidade assintótica espacial e temporal de $O(N^D)$, mas reduz drasticamente as operações constantes internas \cite{barna2024}.
    \item \textbf{Método Integral:} Cria uma imagem integral prévia, reduzindo a complexidade de busca da caixa. Possui tempo e espaço $O(N^D)$, operando em tempo quase idêntico independentemente do tamanho da caixa testada \cite{barna2024}.
    \item \textbf{Método \textit{Trace}:} Baseia-se nos pixels de objeto em vez do tamanho da caixa. Sua complexidade de tempo cai para $O(C)$, onde $C$ é a contagem total de pixels do objeto \cite{barna2024}. Para dados extremamente esparsos, é o algoritmo assintoticamente mais eficiente.
    \item \textbf{Método de Convolução:} Utiliza a Transformada Rápida de Fourier (FFT), obtendo uma complexidade de tempo de $O(N^D \log_2 N)$ e espaço equivalente \cite{barna2024}, ideal quando o formato da "caixa" (ex. janelas circulares) não favorece a imagem integral \cite{barna2024}.
\end{itemize}

Em suma, enquanto o \textit{Box Counting} pode ser resolvido com buscas e ordenações eficientes em tempo log-linear $O(N \log N)$ focadas nos pontos de dados \cite{liebovitch1989}, a exploração estocástica densa fornecida pelo \textit{Gliding Box} exige algoritmos dinâmicos otimizados (como imagens integrais ou abordagem por rastreamento) para evitar que o custo atinja crescimentos polinomiais inviáveis na contagem iterativa bruta \cite{barna2024}.

\begin{algorithm}[htpb]
\caption{Box Counting Otimizado por Ordenação}
\label{alg:box_counting}
\begin{algorithmic}[1] % O [1] numera todas as linhas
    \Require Lista de $N$ pontos, número de resoluções $k$
    \Ensure Dimensão Fractal $D$
    
    \Statex % Linha em branco para separar
    \For{$m = k$ \textbf{até} $0$ \textbf{passo} $-1$}
        \State $MASCARA \gets \text{cria máscara binária para a escala } m$
        
        \For{$i = 1$ \textbf{até} $N$}
            \Comment{Mapeia o ponto para uma caixa específica}
            \State $Hash\_Caixa[i] \gets Coordenada[i] \text{ \textbf{AND} } MASCARA$
        \EndFor
        
        \State \Call{Ordenar}{$Hash\_Caixa$} \Comment{Ex: Quicksort $\mathcal{O}(N \log N)$}
        
        \State $Caixas\_Ocupadas \gets 1$
        \For{$i = 2$ \textbf{até} $N$}
            \If{$Hash\_Caixa[i] \neq Hash\_Caixa[i-1]$}
                \State $Caixas\_Ocupadas \gets Caixas\_Ocupadas + 1$
            \EndIf
        \EndFor
        
        \State \Call{Salvar}{$Caixas\_Ocupadas, \text{tamanho da caixa } 2^m$}
    \EndFor
    \Statex
    \State \Return \Call{RegressãoLinear}{\text{pontos salvos}}
\end{algorithmic}
\end{algorithm}

\begin{algorithm}[htpb]
\caption{Gliding Box Otimizado (Método da Imagem Integral)}
\label{alg:gliding_box}
\begin{algorithmic}[1]
    \Require Imagem $S$ de dimensões $N \times N$, tamanho da caixa $\epsilon$
    \Ensure Lacunaridade $\Lambda(\epsilon)$ para a escala fornecida
    
    \Statex \Comment{Passo 1: Construção da Imagem Integral $I$ -- Tempo $\mathcal{O}(N^2)$}
    \State Criar matriz $I$ de dimensões $(N+1) \times (N+1)$ preenchida com zeros
    \For{$i = 1$ \textbf{até} $N$}
        \For{$j = 1$ \textbf{até} $N$}
            \State $I[i][j] \gets S[i-1][j-1] + I[i-1][j] + I[i][j-1] - I[i-1][j-1]$
        \EndFor
    \EndFor
    
    \Statex \Comment{Passo 2: Cálculo Rápido das Massas -- Tempo $\mathcal{O}(N^2)$ independente de $\epsilon$}
    \State $Tamanho\_M \gets N - \epsilon + 1$
    \State Criar matriz de massas $M$ de dimensões $Tamanho\_M \times Tamanho\_M$
    \For{$i = 0$ \textbf{até} $Tamanho\_M - 1$}
        \For{$j = 0$ \textbf{até} $Tamanho\_M - 1$}
            \State \Comment{Acessa apenas os 4 cantos da matriz integral}
            \State $M[i][j] \gets I[i+\epsilon][j+\epsilon] - I[i][j+\epsilon] - I[i+\epsilon][j] + I[i][j]$
        \EndFor
    \EndFor
    
    \Statex \Comment{Passo 3: Cálculo da Lacunaridade (Equação de Tolle)}
    \State $SomaM \gets 0$
    \State $SomaM2 \gets 0$
    \For{\textbf{cada} massa $m$ \textbf{na matriz} $M$}
        \State $SomaM \gets SomaM + m$
        \State $SomaM2 \gets SomaM2 + (m \times m)$
    \EndFor
    \State $B \gets Tamanho\_M \times Tamanho\_M$ \Comment{Número total de caixas geradas}
    \State $\Lambda(\epsilon) \gets \frac{B \times SomaM2}{SomaM \times SomaM}$
    
    \Statex
    \State \Return $\Lambda(\epsilon)$
\end{algorithmic}
\end{algorithm}

% --- Ambiente thebibliography inserido diretamente no tex ---
\begin{thebibliography}{10}

\bibitem{barna2024}
Barna, B., Kovács, H., \& Erdélyi, M. (2024).
\newblock Methods for calculating gliding-box lacunarity efficiently on large datasets.
\newblock \emph{Pattern Analysis and Applications}, 27(113).

\bibitem{gonzato1998}
Gonzato, G. (1998).
\newblock A practical implementation of the box counting algorithm.
\newblock \emph{Computers \& Geosciences}, 24(1), 95--100.

\bibitem{liebovitch1989}
Liebovitch, L. S., \& Toth, T. (1989).
\newblock A fast algorithm to determine fractal dimensions by box counting.
\newblock \emph{Physics Letters A}, 141(8-9), 386--390.

\bibitem{saa2007}
Saa, A., Gascó, G., Grau, J. B., Antón, J. M., \& Tarquis, A. M. (2007).
\newblock Comparison of gliding box and box-counting methods in river network analysis.
\newblock \emph{Nonlinear Processes in Geophysics}, 14, 603--613.

\end{thebibliography}

\end{document}